\section{Services analytiques F1 à F5}
Les paramètres sont fournis à chaque appel ; les réponses JSON sont calculées
sur les données stockées et renvoient les effectifs et sources. Les exemples
numériques ci-dessous proviennent des réponses archivées du prototype sur janvier
complet, à l'exception du contrôle DST sur des lignes de mars \cite{mesures}.
Ils ne sont pas des résultats recalculés sur les six mois de la version intégrée.

\subsection{F1 : activité au départ ou à l'arrivée}
\code{GET /api/activity} reçoit une zone, \code{pickup} ou \code{dropoff} et un
intervalle. Il renvoie les comptes avant/après qualité, les répartitions par heure,
jour de semaine, date et source. Le filtre utilise la zone et l'horodatage du
sens demandé : une arrivée de février peut provenir d'un départ de janvier.
Après \code{$match} et \code{$project}, \code{$facet} calcule les différentes
répartitions à partir des mêmes candidats ; chaque branche utilise
\code{$group} puis \code{$sort}. Les comptes par motif passent par
\code{$unwind}, qui produit une ligne par flag. Ils peuvent se chevaucher.
Pour les départs de la zone 181 en janvier, le contrôle Parquet archivé donne
983 candidats, 977 retenus et six exclus avec activity. Les distributions sont
cumulées : ce ne sont pas des moyennes corrigées du nombre de lundis ou d'heures.

\subsection{F2 : flux et directions opposées}
\code{GET /api/flows} accepte une zone, un ensemble de zones ou un borough et
une période. Il classe les destinations et paires, compare les retours et expose
les séries hebdomadaires. Le filtre \code{$or} retient les départs \emph{ou}
arrivées dans l'aire pour disposer des deux sens, sans doubler les trajets internes.
Les groupes MongoDB sont définis par origine, destination, semaine et source.
Java combine les petits comptes pour les parts, séries et tris.

Les seuils par défaut demandent 20 trajets pour afficher un lien. Un lien régulier
atteint cinq trajets dans au moins 75\,\% des semaines complètes ; sinon, une
semaine représentant au moins 50\,\% du volume de référence peut le qualifier
de pic. Moins de trois semaines complètes conduit à un historique insuffisant.
Les seuils sont modifiables et les semaines partielles ne servent pas à cette
classification. En janvier, Brooklyn représente 25\,258 candidats et
24\,860 départs retenus par flows. La paire 181 $\to$ 89 compte 48 trajets retenus,
contre 18 en sens inverse ; elle est classée régulière. Le petit retour n'est
pas transformé en flux majeur par une simple variation relative.

\subsection{F3 : activité inhabituelle et couverture}
\code{GET /api/unusual-activity} compare une zone/date/créneau aux mêmes jours
de semaine d'une période antérieure explicite. Il expose les jours retenus,
la médiane $m$, la dispersion $\mathrm{MAD}=\operatorname{mediane}|x_i-m|$ et
les écarts. Les défauts demandent trois références, $m\geq20$, un écart absolu
au moins égal à 20, un écart relatif d'au moins 50\,\% et un score
$|x-m|/(1{,}4826\,\mathrm{MAD})\geq3$. Si MAD vaut zéro, le score est absent,
un avertissement est ajouté et les critères absolu/relatif restent applicables.
C'est un critère descriptif, pas une p-value.

Les jours fériés fédéraux 2024 et les journées entières de changement d'heure
sont exclus des comparaisons usuelles. Le calendrier est embarqué avec sa source
OPM \cite{calendrier}. Un second pipeline compte la ville par jour et créneau :
une absence de mois source, de lignes ville, ou une forte variation générale
peut limiter l'interprétation d'une baisse locale. Ce diagnostic ne prouve pas
une lacune de collecte. Le prototype signale, par exemple, 55 départs retenus
pour la zone 181 le 23 janvier contre une médiane de 20 sur ses références.
La faible profondeur de cet historique et les jours exclus doivent accompagner
cette hausse, sans en déduire une cause.

L'intégration a optimisé le pipeline ville : sur une journée entière, il ne calcule
plus l'heure ; sur un créneau, il extrait directement les deux caractères de
l'heure ISO avec \code{$substrBytes} et \code{$toInt}. Le délai maximal des
agrégations F3 est passé à 300 s. Les temps initiaux ne mesurent pas ces optimisations.

\subsection{F4 : caractéristiques des trajets}
\code{GET /api/journeys} reçoit deux zones et deux périodes non chevauchantes.
Il renvoie médiane, quartiles, p90, moyenne, écart-type, min/max et effectifs pour
la durée et la distance. Les moins fiables sont décrits séparément ; les
montants négatifs et doublons restent hors des statistiques. Les valeurs hors
bornes $[Q_1-1{,}5\,IQR,Q_3+1{,}5\,IQR]$ sont comptées mais pas supprimées à nouveau.

MongoDB filtre la paire et projette durée stockée, distance, calendrier, source
et flags. Un curseur fournit les lignes par lots de 2\,000. Java calcule les
quantiles exacts par interpolation à l'indice $(n-1)p$. Le choix évite un quantile
approximatif ou un grand tableau BSON, au prix d'un transfert et d'une mémoire
proportionnels aux candidats. La limite est de 200\,000 lignes, avec erreur 422
explicite au-delà.

Pour tenir compte de trajets différents, une seconde comparaison utilise les
strates semaine/week-end, créneaux de six heures et classes de distance. Les
moyennes de durée des strates suffisamment représentées sont pondérées par les
mêmes poids de référence ; les fractions couvertes sont affichées. Ce n'est pas
une médiane standardisée ni une élimination de tous les facteurs de confusion.
Pour 132 $\to$ 161, la médiane passe de 2\,678 à 2\,719 s entre les
1--14 et 15--31 janvier, avec 1\,070 et 1\,359 retenus. L'écart médian est
41 s ; la différence des moyennes standardisées vaut 49,76 s. Ces indicateurs
ne suffisent pas à conclure à davantage de congestion.

\subsection{F5 : changement de part face à la variation habituelle}
\code{GET /api/flow-changes} compare des périodes de même durée, lundi à lundi,
et utilise une période historique antérieure au courant. Le dénominateur $N$
est l'ensemble des départs retenus de l'aire sélectionnée, pas seulement les
flux du classement. Pour un flux, $p_A=n_A/N_A$ et $p_B=n_B/N_B$.
La différence $p_B-p_A$ est comparée au p90 des variations absolues entre semaines
historiques adjacentes, propres à ce flux. Le ratio est
\[
R=\frac{|p_B-p_A|}{\max\{q_{0,9}(|p_w-p_{w-1}|),\,\varepsilon\}}.
\]
Les défauts demandent 20 trajets dans au moins une période, 20 dans l'historique,
trois semaines utilisables, un dénominateur hebdomadaire d'au moins 100,
$\varepsilon=0{,}001$ et $R\geq2$. Les flux retenus sont classés par écart absolu
de part, puis par identifiants ; les autres portent leur motif de non-qualification.
Un dénominateur nul donne une part absente, pas un zéro concluant. Le plancher
empêche un score infini sur un historique constant ; il ne rend pas le résultat
statistiquement significatif.

La décomposition $n_B-n_A=(p_A N_B-n_A)+(n_B-p_A N_B)$ distingue l'effet du
volume global et le résidu du flux. Entre les semaines du 15 et du 22 janvier,
les départs retenus de la zone 132 passent de 30\,263 à 26\,739. Pour
132 $\to$ 123, les comptes passent pourtant de 29 à 146, soit de 0,0958\,\% à
0,5460\,\% : +0,4502 point. L'attendu au seul effet volume est 25,62 ; le résidu
est +120,38. La baseline du 1er au 22 janvier ne fournit que deux différences
hebdomadaires, une limite importante. L'intégration construit désormais les
preuves par source après le classement, seulement pour les flux renvoyés, et
porte le délai maximal de l'agrégation à 300 s, sans changer ce principe de calcul.
